\section{The Church Was a Campaign, Not a Date}

By 1872 the community had moved its ambition from reuse to monument. The cornerstone of the present church at Fifth and Mitchell was laid in July. Parish accounts remembered a thirty-dollar assessment on each family---approximately a month of wages for some workers---as the financial base for borrowing and construction. That sum made the church a collective risk: families with little surplus pledged not merely for shelter but for a building that would announce permanence.

Leonard A. Schmidtner designed the church. Later sources identify him with the Polish name Kowalski and a noble title; the documentary basis for every biographical detail remains uneven. His building is more securely read than his legend. Twin towers rise from a cream-brick body; round-arched openings and domes evoke Central European Renaissance and Baroque churches while the preservation inventory classifies the exterior as Romanesque Revival. Later Gothic and Romanizing alterations further complicate any single style label. The church is not a transplanted Polish prototype. It is a Milwaukee composition made from local brick, transatlantic memory, European training, and subsequent campaigns.

A portion was dedicated in 1873 and entered use. Parish histories often call that completion. The federal Historic American Buildings Survey and the Wisconsin Architecture and History Inventory preserve a more revealing account: major work continued or resumed from 1884 through 1894. The sanctuary was built out, the roof received copper, an organ was installed, and the interior was painted. Henry Messmer designed a substantial 1885 renovation valued in the state record at \$17,000. The church people knew at the end of the nineteenth century was therefore not simply the church opened in 1873.

This staged construction explains apparent contradictions in the record. An 1873 dedication can be real, an 1875 photograph can show a recognizable twin-towered church, and an 1880s campaign can still be principal construction. Immigrant parishes often built as money and numbers allowed. Use began before aspiration was exhausted. Saint Stanislaus grew around worship already taking place inside it.

The fabric also preserved the economy of its makers. Cream City brick placed the church within Milwaukee's material identity. Copper domes and clocks made it legible at a distance. Four great bells installed in 1894 carried donor names in their bronze, turning gifts into sound over the neighborhood. The largest was reported at about 5,000 pounds; together the four were said to weigh roughly ten tons. Bells did more than ornament towers. They measured days, announced death and feast, summoned worship, and let an immigrant institution occupy civic air.

\statusframe{Cornerstone, July 1872; partial dedication, 1873; principal later work, 1884--94; Messmer program, 1885; bells, 1894.}{The present church reached recognizable and usable form early, then acquired essential sanctuary, roof, organ, interior, and tower elements through later campaigns.}{The checked sources do not support one all-purpose ``completion date.'' Costs beyond the 1885 state-record figure remain retrospective or unverified.}
